\scanpage{47}
\mspage{46}
\mspage{77}
\margin{n°382}

\S\ 7. \uuline{Jets. Opérateurs différentiels}.

A partir de maintenant, on adopte la convention suivante : si k est un entier, et si r désigne l'un des deux symboles $\infty$ et $\omega$, on pose $r-k=r$. Ainsi les dérivées d'ordre k d'une fonction de classe $C^r$ dans un ouvert de $K^n$ sont de classe $C^{r-k}$ dans cet ouvert, que r soit fini ou non.

7.1 -- \uline{Jets d'applications}.

7.1.1. Soient X et Y deux variétés de classe $C^r$, et soient $x\in X$, $y\in Y$ ; soient f et g deux morphismes d'un voisinage de x à valeurs dans Y, tels que $f(x)=g(x)=y$ ; soit k un entier positif $\leq r$, ou éventuellement le symbole $\infty$ si $r\geq\infty$. On dit que f et g ont un \uline{contact d'ordre} $\geq k$ \uline{en} x s'il existe une carte $(U,\varphi,p)$ de X en x et une carte $(V,\psi,q)$ de Y en y telles que $\psi\circ f\circ\varphi^{-1}$ et $\psi\circ g\circ\varphi^{-1}$ aient un contact d'ordre $\geq k$ au point $\varphi(x)\in K^p$ (cf. n° 1.1.1, 2.5.4). Cette propriété est alors vérifiée quel que soit le choix des cartes $(U,\varphi,p)$ et $(V,\psi,q)$ et définit une relation d'équivalence entre germes de morphismes. Une classe d'équivalence pour cette relation est appelée un \uline{jet d'ordre} k \uline{de morphisme} de X \uline{dans} Y, \uline{de source} x et \uline{de but} y ; la classe du morphisme f est notée $j_x^k(f)$ et se nomme le \uline{jet d'ordre} k de f en x, ou \uline{la partie principale d'ordre} k \uline{de} f en x. L'ensemble des jets d'ordre k de X dans Y (resp. de source x, resp. de but y) est noté $J^k(X,Y)$ (resp. $J_x^k(X,Y)$, resp. $J^k(X,Y)_y$), et l'on pose :
\[
J_x^k(X,Y)_y=J_x^k(X,Y)\cap J^k(X,Y)_y.
\]

7.1.2. Si $Y=K$, $J^k(X,Y)$ est l'ensemble des parties principales d'ordre k de fonctions morphiques sur X. On écrira en général $P^k(X)$, $P_x^k(X)$ et $P_x^k(f)$ pour $J^k(X,K)$, $J_x^k(X,K)$ et $j_x^k(f)$.

Soit $C_x^r(X)$ l'ensemble des germes en x de fonctions morphiques sur un voisinage de x ; c'est une K-algèbre. Il existe sur $P_x^k(X)$ une structure de K-algèbre et une seule telle que l'application
